% Time-stamp: <2026-09-11 16:36:05 komori>
% Copyright (c) 2021 Yasushi Komori
% All rights reserved.

  
%%%%%%%%%%%%%%%%%%%%%%%%%%%%%%%%%%%%%%%%
%%
%% \brl@plot
%%
%% input: 
%%
%% \brl@pgf@minxaty_\brl@pgf@curry : minx at y
%% \brl@pgf@maxxaty_\brl@pgf@curry : maxx at y
%% \brl@pgf@size_\brl@pgf@currx_\brl@pgf@curry : size
%%
%% output: edel form

\def\brl@plot@papersize{B5}
\def\brl@plot@maxlinex_single{480}% 15*32
\def\brl@plot@maxlinex_double{960}% 15*64
\def\brl@plot@maxliney{684}%  (the max of edel : the max of ESA721 is 726)
\def\brl@plot@headsep{0}% 
\def\brl@plot@leftsidemargin{0}%

\def\brl@plot@braillewidth{15}%
\def\brl@plot@brailleheight{31}% important notice : the height of a char of ESA721 is 33.
\def\brl@plot@braillebaseline{22}%


\newwrite\brl@plot@file% for the right page
\newwrite\brl@plot@rfile% for the right page

\newtoks\brl@plot@rtoks% for file of the right page
\def\brl@plot_appendrtoks#1{\brl@plot@rtoks\expandafter{\the\brl@plot@rtoks#1}}

\def\brl@plot{%
  %
  %
  \brl@plot_init%
  \the\brl@pgf@invoketoks%
  % 
  % treat nodes
  %
  \let\brl@bc\brl@bc_plot% make empty box in do_show
  \let\brl@ebc\brl@ebc_plot%
  \let\brl@dbc\brl@dbc_plot% either ebc or dbc works depending on option@hires.
  % \brl@do = \brl@do_show
  \def\brl@pgf@pointsize{2}% the size of dots for letters is 2.
  \let\brl@char@cmd_direct\brl@@asis% may include \label
  \def\brl@pgf@radius{1}%
  \global\setbox\brl@plot@nodebox\hbox{\the\brl@pgf@nodetoks}%
  % 
  \brl@env_plot_grid% make grid. This should be here. after making nodeboxes and before making header. this gives the correct minx, maxx, miny, maxy of the picture.    
  % 
  % treat header
  % 
  % 
  \global\setbox\brl@plot@headerbox\hbox{%
    \let\brl@pgf_orient\brl@pgf_orient_portrait% force portrait
    % 
    % 
    \brl@pgf_hbox{\numexpr\brl@pgf@minx+356\relax}{\numexpr\brl@pgf@maxy+17\relax}{\brl@plot@page}% 15 * 24 - 4 spaces = 356, 31 * 1 = 31 (1 lineskip) 31-14 (height)     
    % \@namedef{brl@bc_plot@x0}{4}
    % \@namedef{brl@bc_plot@y0}{14}
    % 
    % 
  }%
  % 
  % 
  % 
  \edef\brl@plot@boundingboxx{\the\numexpr\brl@plot@maxlinex+\brl@pgf@minx-\brl@plot@leftsidemargin\relax}%
  \edef\brl@plot@boundingboxy{\the\numexpr\brl@pgf@maxy+\brl@plot@headsep-\brl@plot@maxliney\relax}%
  % 
  \brl@plot@liney\brl@pgf@maxy\relax%
  % 
  % 
  % 
  % 
  % 
  % 
  \brl@@pdf_bcontent%
  \brl@@pdf_invoke{q 0 G 0 g 0 w}%
  \brl@@pdf_invoke{q 1 0 0 1 \the\numexpr\brl@plot@leftsidemargin-\brl@pgf@minx\relax\brl@@space%
    \the\numexpr-\brl@pgf@maxy-\brl@plot@headsep\relax\brl@@space cm}%
  % 
  % init of writeps
  % 
  \def\brl@plot@currsize{0}%
  \edef\brl@tmp{{EDEL\brl@plot@papersize,\brl@plot@headsep,%
      \the\numexpr\brl@pgf@maxy-\brl@pgf@miny+\brl@plot@headsep\relax%
    }}%
  \brl@plot@toks\brl@tmp% header of EDEL
  % 
  % 
  \ifbrl@figure@spread%
    \def\brl@plot@rcurrsize{0}%
    \brl@plot@rtoks\brl@tmp% rheader of EDEL
    \let\brl@plot_writepts_subroutine\brl@plot_writepts_double%
    \immediate\openout\brl@plot@file=\jobname-\brl@figure@figurenum_plot-l.edl\relax%
    \immediate\openout\brl@plot@rfile=\jobname-\brl@figure@figurenum_plot.edl\relax%
  \else%
    \immediate\openout\brl@plot@file=\jobname-\brl@figure@figurenum_plot.edl\relax%
  \fi%
  % 
  % 
  \brl@plot_writepts%
  \immediate\write\brl@plot@file{\the\brl@plot@toks}% treat the rest
  \immediate\closeout\brl@plot@file%
  % 
  \brl@@pdf_invoke{Q 1 0 0 1 -0.5 -0.5 cm}% i.e. -0.5 -- 479.5 etc.
  % 
  {\brl@orig@color{\brl@char@warncolor}%
    \brl@@pdf_invoke{0 5 \brl@plot@maxlinex_single\brl@@space-\the\numexpr\brl@plot@maxliney+5\relax\brl@@space re S}%  +5 means the radius of the top dots of the header
    % 
    % 
    \ifbrl@figure@spread%
      \immediate\write\brl@plot@rfile{\the\brl@plot@rtoks}% treat the rest
      \immediate\closeout\brl@plot@rfile%
      % 
      \brl@@pdf_invoke{\brl@plot@maxlinex_single\brl@@space 5 \brl@plot@maxlinex_single\brl@@space-\the\numexpr\brl@plot@maxliney+5\relax\brl@@space re S}% 
    \fi%
  }%
  \brl@@pdf_invoke{Q}%
  \brl@@pdf_econtent%
  % 
}



\def\brl@plot_writepts{%
  \ifcsname brl@pgf@minxaty_\the\brl@plot@liney\endcsname%
    \edef\brl@plot@maxxaty{\csname brl@pgf@maxxaty_\the\brl@plot@liney\endcsname}%
    \brl@plot@linex\csname brl@pgf@minxaty_\the\brl@plot@liney\endcsname\relax%
    %
    \global\expandafter\let\csname brl@pgf@minxaty_\the\brl@plot@liney\endcsname\@undefined% clear
    %
    \ifnum\brl@plot@boundingboxy<\brl@plot@liney%
      % 
      \begingroup%
        \count@\numexpr\brl@pgf@maxy-\brl@plot@liney+\brl@plot@headsep\relax% upside down
        \brl@plot@liney\count@%
        \divide\count@26\relax%
        \lccode`\!=\numexpr\count@+64\relax%
        \lccode`\?=\numexpr\brl@plot@liney-\count@*26+64\relax%
        \lowercase{\gdef\brl@plot@lineychar{!?}}%
      \endgroup%
      \brl@plot_writepts_loop%
    \else%
      \brl@plot_writepts_outy_loop%
    \fi%
  \fi%
  %
  %
  %
  %
  \ifnum\brl@pgf@miny<\brl@plot@liney%
    \advance\brl@plot@liney\m@ne%
    \expandafter\brl@plot_writepts%
  \fi%
}


\def\brl@plot_writepts_outy_loop{%
  \ifcsname brl@pgf@size_\the\brl@plot@linex _\the\brl@plot@liney\endcsname%
    %
    %
    \brl@@pdf_invoke{q 1 0 0 1 \the\brl@plot@linex\brl@@space\the\brl@plot@liney\brl@@space cm}%
    \brl@plot@circle{\csname brl@pgf@size_\the\brl@plot@linex _\the\brl@plot@liney\endcsname r}% draw circle
    \brl@@pdf_invoke{Q}%
    % 
    %
    \global\expandafter\let\csname brl@pgf@size_\the\brl@plot@linex _\the\brl@plot@liney\endcsname\@undefined% clear
    %
  \fi%
  %
  %
  %
  \ifnum\brl@plot@maxxaty>\brl@plot@linex%
    \advance\brl@plot@linex\@ne%
    \expandafter\brl@plot_writepts_outy_loop%
  \fi%
}


\def\brl@plot_writepts_loop{%
  \ifcsname brl@pgf@size_\the\brl@plot@linex _\the\brl@plot@liney\endcsname%
    % whether out of range or not 
    \brl@@pdf_invoke{q 1 0 0 1 \the\brl@plot@linex\brl@@space\the\brl@plot@liney\brl@@space cm}%
    %
    \edef\brl@tmp{\csname brl@pgf@size_\the\brl@plot@linex _\the\brl@plot@liney\endcsname}% size
    %
    \ifnum\brl@plot@boundingboxx>\brl@plot@linex%
      \brl@plot@circle{\brl@tmp}% draw circle
      \count@\numexpr\brl@plot@linex-\brl@pgf@minx+\brl@plot@leftsidemargin\relax% x pos
      % 
      \brl@plot_writepts_subroutine% use \brl@tmp = size, \count@
    \else%
      \brl@plot@circle{\brl@tmp r}% draw circle
    \fi%
    % 
    \brl@@pdf_invoke{Q}%
    % 
    %
    \global\expandafter\let\csname brl@pgf@size_\the\brl@plot@linex _\the\brl@plot@liney\endcsname\@undefined% clear
    %
  \fi%
  %
  %
  %
  \ifnum\brl@plot@maxxaty>\brl@plot@linex%
    \advance\brl@plot@linex\@ne%
    \expandafter\brl@plot_writepts_loop%
  \fi%
}


\def\brl@plot_writepts_double{%
  \ifnum\brl@plot@maxlinex_single>\count@%
    % 
    % 
    \ifnum\brl@tmp=\brl@plot@currsize\relax% within the single page
    \else% different size
      %
      \let\brl@plot@currsize\brl@tmp%
      \immediate\write\brl@plot@file{\the\brl@plot@toks}%
      \brl@plot@toks\expandafter{\brl@plot@currsize}%
    \fi%
    \begingroup%
      \brl@plot@linex\count@%
      \divide\count@26\relax%
      \lccode`\!=\numexpr\count@+64\relax%
      \lccode`\?=\numexpr\brl@plot@linex-\count@*26+64\relax%
      \lowercase{\xdef\brl@tmp{{!?\brl@plot@lineychar}}}%
    \endgroup%
    %
    \expandafter\brl@plot_appendtoks\brl@tmp%
    % 
  \else%
    % 
    % 
    \ifnum\brl@tmp=\brl@plot@rcurrsize\relax% in the right page
    \else% different size
      %
      \let\brl@plot@rcurrsize\brl@tmp%
      \immediate\write\brl@plot@rfile{\the\brl@plot@rtoks}%
      \brl@plot@rtoks\expandafter{\brl@plot@rcurrsize}%
    \fi%
    \begingroup%
      \advance\count@-\brl@plot@maxlinex_single\relax%
      \brl@plot@linex\count@%
      \divide\count@26\relax%
      \lccode`\!=\numexpr\count@+64\relax%
      \lccode`\?=\numexpr\brl@plot@linex-\count@*26+64\relax%
      \lowercase{\xdef\brl@tmp{{!?\brl@plot@lineychar}}}%
    \endgroup%
    % 
    \expandafter\brl@plot_appendrtoks\brl@tmp%
    %
  \fi%
}

% default
\def\brl@plot_writepts_subroutine{%
  %
  \ifnum\brl@tmp=\brl@plot@currsize\relax% within the single page
  \else% different size
    %
    \let\brl@plot@currsize\brl@tmp%
    \immediate\write\brl@plot@file{\the\brl@plot@toks}%
    \brl@plot@toks\expandafter{\brl@plot@currsize}%
  \fi%
  \begingroup%
    \brl@plot@linex\count@%
    \divide\count@26\relax%
    \lccode`\!=\numexpr\count@+64\relax%
    \lccode`\?=\numexpr\brl@plot@linex-\count@*26+64\relax%
    \lowercase{\xdef\brl@tmp{{!?\brl@plot@lineychar}}}%
  \endgroup%
  \expandafter\brl@plot_appendtoks\brl@tmp%
  % 
}



\ifbrl@lualatex%
  %
\expandafter\newbox\csname brl@plot@circle1_\endcsname
\expandafter\newbox\csname brl@plot@circle1r_\endcsname
\expandafter\newbox\csname brl@plot@circle2_\endcsname
\expandafter\newbox\csname brl@plot@circle2r_\endcsname
\expandafter\newbox\csname brl@plot@circle3_\endcsname
\expandafter\newbox\csname brl@plot@circle3r_\endcsname
\def\brl@plot_resource{%
  %
  \setbox\@ne\hbox{}\wd\@ne\p@\ht\@ne\p@%
  \setbox\z@\hbox{%
    \pdfextension literal{0 g \brl@@cira}%
    \copy\@ne%
  }%
  \saveboxresource type 2 attr{%
/Subtype /Form
/Type /XObject
/FormType 1
/BBox[-1 -1 1 1]}\z@%
  \setbox\z@\hbox{\useboxresource\lastsavedboxresourceindex}%
  \wd\z@\z@%
  \ht\z@\z@%
  \dp\z@\z@%
  \setbox\csname brl@plot@circle1_\endcsname\box\z@%  
  %
  \setbox\z@\hbox{%
    \pdfextension literal{1 0 0 rg \brl@@cira}%
    \copy\@ne%                             
  }%
  \saveboxresource type 2 attr{%
/Subtype /Form
/Type /XObject
/FormType 1
/BBox[-1 -1 1 1]}\z@%
  \setbox\z@\hbox{\useboxresource\lastsavedboxresourceindex}%
  \wd\z@\z@%
  \ht\z@\z@%
  \dp\z@\z@%
  \setbox\csname brl@plot@circle1r_\endcsname\box\z@%  
  %
  \setbox\z@\hbox{%
    \pdfextension literal{0 g \brl@@cirb}%
    \copy\@ne%
  }%
  \saveboxresource type 2 attr{%
/Subtype /Form
/Type /XObject
/FormType 1
/BBox[-2 -2 2 2]}\z@%
  \setbox\z@\hbox{\useboxresource\lastsavedboxresourceindex}%
  \wd\z@\z@%
  \ht\z@\z@%
  \dp\z@\z@%
  \setbox\csname brl@plot@circle2_\endcsname\box\z@%  
  %
  \setbox\z@\hbox{%
    \pdfextension literal{1 0 0 rg \brl@@cirb}%
    \copy\@ne%
  }%
  \saveboxresource type 2 attr{%
/Subtype /Form
/Type /XObject
/FormType 1
/BBox[-2 -2 2 2]}\z@%
  \setbox\z@\hbox{\useboxresource\lastsavedboxresourceindex}%
  \wd\z@\z@%
  \ht\z@\z@%
  \dp\z@\z@%
  \setbox\csname brl@plot@circle2r_\endcsname\box\z@%  
  %
  \setbox\z@\hbox{%
    \pdfextension literal{0 g \brl@@circ}%
    \copy\@ne%
  }%
  \saveboxresource type 2 attr{%
/Subtype /Form
/Type /XObject
/FormType 1
/BBox[-3 -3 3 3]}\z@%
  \setbox\z@\hbox{\useboxresource\lastsavedboxresourceindex}%
  \wd\z@\z@%
  \ht\z@\z@%
  \dp\z@\z@%
  \setbox\csname brl@plot@circle3_\endcsname\box\z@%  
  %
  \setbox\z@\hbox{%
    \pdfextension literal{1 0 0 rg \brl@@circ}%
    \copy\@ne%
  }%
  \saveboxresource type 2 attr{%
/Subtype /Form
/Type /XObject
/FormType 1
/BBox[-3 -3 3 3]}\z@%
  \setbox\z@\hbox{\useboxresource\lastsavedboxresourceindex}%
  \wd\z@\z@%
  \ht\z@\z@%
  \dp\z@\z@%
  \setbox\csname brl@plot@circle3r_\endcsname\box\z@%  
  %
}
\def\brl@plot@circle#1{%
  \copy\csname brl@plot@circle#1_\endcsname% draw circle
}%

\else%

\def\brl@plot_resource{%
  \AtBeginDvi{%
    \special{pdf:bxobj @circle1 bbox -1 -1 1 1}%
    \special{pdf:code 0 g \brl@@cira}%
    \special{pdf:exobj}%
    % 
    \special{pdf:bxobj @circle1r bbox -1 -1 1 1}%
    \special{pdf:code 1 0 0 rg \brl@@cira}%
    \special{pdf:exobj}%
    % 
    \special{pdf:bxobj @circle2 bbox -2 -2 2 2}%
    \special{pdf:code 0 g \brl@@cirb}%
    \special{pdf:exobj}%
    % 
    \special{pdf:bxobj @circle2r bbox -2 -2 2 2}%
    \special{pdf:code 1 0 0 rg \brl@@cirb}%
    \special{pdf:exobj}%
    % 
    \special{pdf:bxobj @circle3 bbox -3 -3 3 3}%
    \special{pdf:code 0 g \brl@@circ}%
    \special{pdf:exobj}%
    % 
    \special{pdf:bxobj @circle3r bbox -3 -3 3 3}%
    \special{pdf:code 1 0 0 rg \brl@@circ}%
    \special{pdf:exobj}%
    % 
  }%
}


\def\brl@plot@circle#1{%
  \special{pdf:uxobj @circle#1}% draw circle
}%

\fi


\endinput


%%% Local Variables:
%%% mode: japanese-latex
%%% backup-each-save-backup : t
%%% TeX-master: t
%%% End:
